% Section 8.7, from lectures/L08/README.md: the objectives, the questions to test yourself, and the
% README's "Next lecture", which in the last lecture is the end of the course.

\needspace{10\baselineskip}
\section{Review}
\label{c8:sec:review}

\subsection*{What you should be able to do now}
After this chapter, you should be able to:
\begin{itemize}
  \item Explain what a finite state machine is, and identify its states, transitions, inputs and
    outputs in a given circuit.
  \item Design a Moore machine by hand, from specification through state diagram and state table to
    Karnaugh-derived equations, and build the result in CircuitVerse.
  \item Model the same machine in VHDL with an enumerated \code{state_t} type and a \code{case}
    process.
  \item Recognize a Mealy machine on sight and say what its output costs and buys, one clock cycle
    of latency against one extra state, and why Moore is the default in this book.
  \item Compose a state machine with previously built subcomponents rather than rewriting them. By
    the capstone that means three earlier chapters' modules around a machine that is the only new
    logic, drawing on the shift and counter idioms of one more.
\end{itemize}

\needspace{6\baselineskip}
\subsection*{Questions to test yourself}
You should be able to explain:
\begin{enumerate}
  \item What makes a circuit a state machine, rather than just a register with some logic around
    it?
  \item In the hand-designed machine, what exactly is \code{X}, and why must it be a one-cycle
    pulse rather than the button's level?
  \item Why is that machine's state encoding a Gray code? What could go wrong with an ordinary
    binary count?
  \item A counter and a shift register are technically state machines too. What distinguishes them
    from the machine designed here?
  \item Which signals is a Moore output allowed to depend on, and a Mealy output? Given a single
    line of VHDL driving an output, how do you tell which you are looking at?
  \item In \secref{c8:a5}'s comparison, why does the Moore version need an extra state, and why does
    its output lag the Mealy version's by exactly one clock cycle?
  \item In \file{fsm_led.vhd}, why does \code{LED_PROCESS} only ever read \code{state}, never
    \code{button_edge_s2} or \code{button_n} directly? What would change if it did?
  \item Why is \code{din} in \code{seq_detect_mealy} not passed through a synchronizer, unlike
    \code{reset_n}?
  \item What does \code{when others} buy you in a state-transition \code{case}, given that
    \code{state_t} has only the values you declared?
  \item In the capstone receiver, why does the timer run at sixteen times the baud rate rather than
    once per bit, and why does sampling in the middle of a bit make the design tolerate a
    transmitter whose clock is slightly different from yours?
  \item The capstone transmitter uses a whole-bit timer and needs no mid-bit anything. What does it
    have that the receiver doesn't, that makes the difference?
\end{enumerate}

\needspace{6\baselineskip}
\subsection*{Where this leaves you}
This is the final chapter of \emph{Digital Design with VHDL}. Over eight chapters the book moved
from gates and Karnaugh maps, through registers, metastability, and the VHDL language, to counters,
shift registers and timers, and now finite state machines, at every step reasoning about a circuit
by hand before writing it in synthesizable VHDL and watching it run on a real FPGA. Every pattern
built along the way, the double-flop synchronizer, the \code{process}/\code{case} coding style, and
now the Moore state machine, was reused again in this chapter's own examples, and carries forward
unchanged into whatever you build next.
